\section{The one-step estimate}\label{ch18:sec:onestep}
Write
\[
 e_n=d(y_n,x_*)
\]
for the distance from the $n$th point to the fixed point. We call $e_n$ the \emph{error} at step~$n$. Keep the two kinds of error apart: $e_n$ measures how far the point $y_n$ is from $x_*$, while the update error $\eps_n$ measures how badly a single update went wrong.

To compare $e_{n+1}$ with $e_n$, insert the exact update $T(y_n)$ as a stop between $y_{n+1}$ and $x_*$:
\begin{align*}
 e_{n+1}
 &=d(y_{n+1},x_*)\\
 &\le d(y_{n+1},T(y_n))+d(T(y_n),x_*)
 &&\text{triangle inequality}\\
 &=d(y_{n+1},T(y_n))+d(T(y_n),T(x_*))
 &&\text{since $T(x_*)=x_*$}\\
 &\le\eps_n+q\,d(y_n,x_*)
 &&\text{update error and contraction inequality}\\
 &=q e_n+\eps_n.
\end{align*}
So for every $n\in\NN$ we have the \emph{one-step estimate}
\[
 e_{n+1}\le qe_n+\eps_n.
\]
Each assumption is used exactly once: the triangle inequality of the metric, the fixed-point equation $T(x_*)=x_*$, the bound on the update error, and the contraction inequality. Nothing else is needed. Habit~3 asks which information an argument actually uses, and here the answer is short. That short list also shows how far the estimate reaches: $X$ need not be a set of numbers, and the update errors need not follow any pattern.

\subsection*{Unfold three steps before writing the general sum}
The one-step estimate compares each error only with the one just before it. To see how an update error made early on affects later steps, apply the estimate repeatedly, substituting each bound into the next. Every substitution multiplies an inequality by $q$, which keeps the direction of the inequality because $q>0$:
\begin{align*}
 e_1&\le qe_0+\eps_0,\\
 e_2&\le qe_1+\eps_1
       \le q(qe_0+\eps_0)+\eps_1
       =q^2e_0+q\eps_0+\eps_1,\\
 e_3&\le qe_2+\eps_2
       \le q(q^2e_0+q\eps_0+\eps_1)+\eps_2
       =q^3e_0+q^2\eps_0+q\eps_1+\eps_2.
\end{align*}
At step three, the starting error $e_0$ has been multiplied by $q$ three times, once for each update. The update error $\eps_0$, made in the first update, has been multiplied by $q$ twice, once for each of the two updates that came after it. The most recent update error $\eps_2$ has not yet been reduced by any later update, so its coefficient is one.

The same counting works at every step. The update error $\eps_j$ is made while producing $y_{j+1}$. After that, the updates that produce $y_{j+2},\ldots,y_n$ follow, and there are $n-(j+1)=n-1-j$ of them. Each multiplies the effect of $\eps_j$ by $q$, so in the bound for $e_n$ the error $\eps_j$ carries the weight $q^{n-1-j}$. As a check, the latest update error has $j=n-1$ and weight $q^0=1$, and the earliest has $j=0$ and weight $q^{n-1}$. With $n=3$ this gives the weights $q^2$, $q$ and $1$ found above. These two endpoint checks fix the limits of the sum in the theorem below.

\subsection*{Two kinds of update errors}
The theorem treats two kinds of update errors separately. The first kind is \emph{bounded}: there is a fixed number $\eta\ge0$ with $\eps_j\le\eta$ for every $j$. Such errors may stay the same size forever, like the constant bias of Section~\ref{ch18:sec:guess}. The second kind \emph{tends to zero}: $\eps_j\to0$ in the sense of Definition~\ref{ch11:def:convergence}, which here means that for every tolerance $\delta>0$, all the update errors from some index onward are smaller than $\delta$.

Errors of the second kind can still add up to a lot. For example, $\eps_j=1/(j+1)$ tends to zero by Section~\ref{ch11:sec:convergence}, but the total $\eps_0+\eps_1+\cdots+\eps_{n-1}=1+\frac12+\cdots+\frac1n$ is a harmonic sum, and Section~\ref{ch11:sec:cauchy} showed that these sums do not converge. We will see that $y_n\to x_*$ all the same. Bounded errors give a bound on the distance to $x_*$; errors that tend to zero give convergence. The two conclusions need different arguments, so the theorem states them separately.

\par\noindent\begin{minipage}{\linewidth}
\begin{theorem}[Iteration with update errors]\label{thm:inexact}
Let $(X,d)$ be a nonempty complete metric space, let $T:X\to X$ be a contraction with factor $q$, where $0<q<1$, and let $x_*$ be its fixed point. Let $y_0,y_1,y_2,\ldots$ be points of $X$ and let $\eps_0,\eps_1,\eps_2,\ldots$ be nonnegative numbers such that $d(y_{n+1},T(y_n))\le\eps_n$ for every $n\in\NN$. Write $e_n=d(y_n,x_*)$.
\begin{enumerate}
\item For every integer $n\ge1$,
\[
 e_n\le q^n e_0+\sum_{j=0}^{n-1}q^{n-1-j}\eps_j.
\]
\item If there is a number $\eta\ge0$ with $\eps_j\le\eta$ for every $j\in\NN$, then for every integer $n\ge1$,
\[
 e_n\le q^n e_0+\eta\,\frac{1-q^n}{1-q}.
\]
\item If $\eps_n\to0$, then $y_n\to x_*$.
\end{enumerate}
\end{theorem}
\end{minipage}\par

Before the proof, here is how to read part~2. The first term, $q^ne_0$, is what remains of the starting error, and it tends to zero as $n$ grows. The second term is the combined effect of all the update errors so far. Since $1-q^n<1$, it is never larger than $\eta/(1-q)$, and it approaches $\eta/(1-q)$ as $n$ grows. So after many steps the error is at most a little more than $\eta/(1-q)$. Section~\ref{ch18:sec:sharp} shows that the factor $1/(1-q)$ cannot be lowered.

\begin{proof}
\textbf{Part 1.} We use induction on $n$. For $n=1$, the sum has the single term $j=0$, whose weight is $q^0=1$, so the claim reads $e_1\le qe_0+\eps_0$. This is the one-step estimate with $n=0$. Now suppose the claim holds for some $n\ge1$. Apply the one-step estimate, and then the induction hypothesis multiplied by the positive number $q$:
\begin{align*}
 e_{n+1}
 &\le qe_n+\eps_n\\
 &\le q\left(q^n e_0+\sum_{j=0}^{n-1}q^{n-1-j}\eps_j\right)+\eps_n\\
 &=q^{n+1}e_0+\sum_{j=0}^{n-1}q^{n-j}\eps_j+q^0\eps_n\\
 &=q^{n+1}e_0+\sum_{j=0}^{n}q^{n-j}\eps_j.
\end{align*}
In the third line, $q\cdot q^{n-1-j}=q^{n-j}$, and the last term $\eps_n$ is written as $q^0\eps_n$ so that it becomes the term $j=n$ of the new sum. Since $n-j=(n+1)-1-j$, the last line is the claim with $n+1$ in place of $n$. By induction, the claim holds for every $n\ge1$.

\textbf{Part 2.} Suppose $\eps_j\le\eta$ for every $j$. Every weight $q^{n-1-j}$ is positive, so replacing each $\eps_j$ by the number $\eta$, which is at least as large, cannot make the sum smaller. Then put $k=n-1-j$. As $j$ runs from $0$ up to $n-1$, the exponent $k$ runs from $n-1$ down to $0$, so the same powers of $q$ appear in the opposite order:
\[
 \sum_{j=0}^{n-1}q^{n-1-j}\eps_j
 \le\eta\sum_{j=0}^{n-1}q^{n-1-j}
 =\eta\sum_{k=0}^{n-1}q^k
 =\eta\,\frac{1-q^n}{1-q}.
\]
The last equality is the finite geometric sum of Section~\ref{ch11:sec:geometric}, with $r=n$. Together with part~1, this proves part~2.

\textbf{Part 3: the plan.} Now suppose $\eps_n\to0$. We must show that for every tolerance $\delta>0$ there is a threshold $N$ such that $e_n<\delta$ for every $n\ge N$.

The sum in part~1 mixes two kinds of update errors. The early ones may be large, but there are only finitely many of them, and every later update multiplies their effect by another factor $q$. The late ones may be very many, but each of them is small, and their weights $q^{n-1-j}$ add up to less than $1/(1-q)$. This suggests splitting the sum at some index $J$. Before $J$ we rely on the shrinking powers of $q$; from $J$ onward we rely on the update errors being small.

The order of the choices matters. First the tolerance $\delta$ is given. Then we choose $J$, depending on $\delta$, so that every update error from index $J$ onward is small. Only after $J$ is fixed do we choose the threshold $N$, large enough to shrink the early part. The index $J$ must not change with $n$: the early part has to be one fixed collection of finitely many terms, so that a single threshold $N$ can make it small for every $n\ge N$ at once.

\textbf{Part 3: choosing $J$.} Let $\delta>0$. The number $(1-q)\delta/2$ is positive, and $\eps_j\to0$, so the definition of convergence, applied with this tolerance and with the distance $|\eps_j-0|=\eps_j$, gives an index $J$ such that
\[
 \eps_j<\frac{(1-q)\delta}{2}\qquad\text{for every }j\ge J.
\]
A larger index works as well, so we may take $J\ge1$. For every $n>J$, split the sum of part~1 into the terms with $j<J$ and the terms with $j\ge J$:
\[
 e_n\le\left(q^ne_0+\sum_{j=0}^{J-1}q^{n-1-j}\eps_j\right)+\sum_{j=J}^{n-1}q^{n-1-j}\eps_j.
\]

\textbf{Part 3: the early part.} Written out, the expression in parentheses is
\[
 q^ne_0+q^{n-1}\eps_0+q^{n-2}\eps_1+\cdots+q^{n-J}\eps_{J-1}.
\]
The smallest exponent among these terms is $n-J$. Factoring out $q^{n-J}$ leaves
\[
 C_J=q^Je_0+q^{J-1}\eps_0+q^{J-2}\eps_1+\cdots+\eps_{J-1},
\]
so the early part equals $q^{n-J}C_J$. The number $C_J$ is a sum of finitely many fixed nonnegative numbers. It depends on $J$, but not on $n$.

\textbf{Part 3: the late part.} Every $\eps_j$ with $j\ge J$ is less than $(1-q)\delta/2$, and the weights are positive, so
\[
 \sum_{j=J}^{n-1}q^{n-1-j}\eps_j
 <\frac{(1-q)\delta}{2}\sum_{j=J}^{n-1}q^{n-1-j}
 =\frac{(1-q)\delta}{2}\sum_{k=0}^{n-1-J}q^k.
\]
Here we substituted $k=n-1-j$ again: as $j$ runs from $J$ up to $n-1$, the exponent $k$ runs from $n-1-J$ down to $0$. By the finite geometric sum, $\sum_{k=0}^{n-1-J}q^k=(1-q^{n-J})/(1-q)$, which is less than $1/(1-q)$. Hence the late part is less than
\[
 \frac{(1-q)\delta}{2}\cdot\frac1{1-q}=\frac\delta2.
\]
This explains the tolerance we asked of the update errors. The factor $1-q$ in it cancels the factor $1/(1-q)$ that comes from adding up the weights, and what is left is half of $\delta$. The other half is kept for the early part.

\textbf{Part 3: choosing $N$.} So far, for every $n>J$,
\[
 e_n<q^{n-J}C_J+\frac\delta2.
\]
If $C_J=0$, this says $e_n<\delta/2<\delta$ for every $n>J$, and $N=J+1$ will do. If $C_J>0$, recall from Section~\ref{ch11:sec:geometric} that $q^m\to0$. With the tolerance $\delta/(2C_J)$, this gives a threshold $M$ such that $q^m<\delta/(2C_J)$ for every $m\ge M$, and we may take $M\ge1$. Put $N=J+M$. For every $n\ge N$ we have $n>J$ and $n-J\ge M$, so $q^{n-J}C_J<\delta/2$, and therefore
\[
 e_n<\frac\delta2+\frac\delta2=\delta.
\]
In both cases $d(y_n,x_*)=e_n<\delta$ for every $n\ge N$. Since $\delta>0$ was arbitrary, this proves $y_n\to x_*$.
\end{proof}

Splitting off finitely many early terms is a technique worth remembering. The early update errors may be large, but there are only finitely many of them, and their effect is eventually damped by the powers of $q$. The late ones are controlled all together, because from index $J$ onward every one of them is small. It is the same idea as splitting a tolerance into two halves in Section~\ref{ch11:sec:cauchy}, now combined with a split of the sum.

Notice also what the proof did not need. It never required the total $\eps_0+\eps_1+\cdots+\eps_n$ to stay bounded; for $\eps_j=1/(j+1)$ it does not, and still $y_n\to x_*$. What the proof controls is the weighted sum, in which each update error is multiplied by a power of $q$ that becomes smaller the older the error is.

\pauseq{Part~3 of Theorem~\ref{thm:inexact} assumes only that $\eps_n\to0$. Does it tell us how many steps are needed to get within $0.001$ of $x_*$?}{No. The proof produces a threshold $N$ for each tolerance, but $N$ depends on the index $J$, and $J$ depends on how quickly the update errors shrink, which part~3 does not specify. If the update errors shrink very slowly, the points may approach $x_*$ very slowly too. When more is known about the update errors, more can be said: Exercise~\ref{ex:18.4} gives an explicit bound when they shrink geometrically.}
